\documentclass{beamer}

\usepackage[T1]{fontenc}
\usepackage[utf8]{inputenc}
\usepackage{hyperref}

\definecolor{plrgprimary}{HTML}{1F2A44}
\definecolor{plrgaccent}{HTML}{3730A3}
\definecolor{plrgfunctional}{HTML}{0F766E}
\definecolor{plrgbackground}{HTML}{F8FAFC}
\definecolor{plrgtext}{HTML}{111827}

\usefonttheme{professionalfonts}

\setbeamercolor{background canvas}{bg=plrgbackground}
\setbeamercolor{title}{fg=plrgprimary}
\setbeamercolor{frametitle}{fg=plrgprimary}
\setbeamercolor{subtitle}{fg=plrgprimary}
\setbeamercolor{structure}{fg=plrgaccent}
\setbeamercolor{normal text}{fg=plrgtext}

\setbeamerfont{title}{series=\bfseries}
\setbeamerfont{frametitle}{series=\bfseries}

\setbeamertemplate{navigation symbols}{}
\setbeamertemplate{footline}[frame number]

\hypersetup{
  colorlinks=true,
  urlcolor=plrgfunctional,
}

\title{Proof Assistants}
\subtitle{Introduction to Program Verification --- Part 1}
\author{Jaemin Hong}
\date{}

\begin{document}

\begin{frame}
  \titlepage
\end{frame}

\begin{frame}{Install Lean}
  \begin{itemize}
    \item 1. Install VS Code
    \item 2. Install the Lean 4 extension
    \item Refer to \url{https://lean-lang.org/install}
  \end{itemize}
\end{frame}

\begin{frame}{Program Verification}
  \begin{itemize}
    \item Proving a property of a program
    \item Some interesting properties
      \begin{itemize}
        \item Absence of memory safety violations
        \item Information flow security
        \item (Non-)termination
        \item Type soundness
        \item Compiler correctness
        \item This program computes the Fibonacci number of its input
      \end{itemize}
  \end{itemize}
\end{frame}

\begin{frame}{Comparison with Other Approaches}
  \begin{itemize}
    \item Testing
      \begin{itemize}
        \item Can prove the presence of bugs, but not their absence
        \item Typically requires an executable oracle
      \end{itemize}
    \item (Sound) static analysis
      \begin{itemize}
        \item Suffers from false positives
        \item Each analyzer targets a specific class of properties
      \end{itemize}
  \end{itemize}
\end{frame}

\begin{frame}{Validation of Proofs}
  \begin{itemize}
    \item How can we trust a proof consisting of thousands of lines?
  \end{itemize}
\end{frame}

\begin{frame}{Proof Assistants}
  \begin{quote}
    \textbf{Formal proof assistants} are pieces of software designed to help their
    users carry out computer-checked proofs. We usually call them \textbf{proof
    assistants}, or \textbf{interactive theorem provers}, but a frustrated student
    coined the phrase ``proof-preventing beasts,'' and dictation software
    occasionally misunderstands ``theorem prover'' as ``fear improver.''
  \end{quote}

  (Baanen et al., The Hitchhiker's Guide to Logical Verification\footnote{\url{https://github.com/lean-forward/logical_verification_2025}})
\end{frame}

\begin{frame}{Formal Proofs}
  \begin{quote}
    A \textbf{formal proof} is a logical argument expressed in a logical formalism.
    In this context, ``formal'' means ``logical'' or ``logic-based.''
  \end{quote}

  \begin{quote}
    An \textbf{informal proof} is what a mathematician would normally call a proof.
    The level of detail can vary a lot, and phrases such as ``it is obvious
    that,'' ``clearly,'' and ``without loss of generality'' move some of the
    proof burden onto the reader.
  \end{quote}

  (Baanen et al., The Hitchhiker's Guide to Logical Verification)
\end{frame}

\begin{frame}{Different Proof Assistants}
  \begin{itemize}
    \item Grouped by their logical foundation
      \begin{itemize}
        \item Dependent type theory: Agda, Lean, Rocq (formerly known as Coq)
        \item Simple type theory: Isabelle/HOL
        \item Set theory: Isabelle/ZF
      \end{itemize}
  \end{itemize}
\end{frame}

\begin{frame}{Formal Verification in the Real World}
  \begin{itemize}
    \item CompCert, CakeML: verified compilers
    \item seL4, CertiKOS: verified kernels
  \end{itemize}
\end{frame}

\begin{frame}{Formal Proofs in the AI Era}
  \begin{quote}
    We are sharing a selection of ten results, each of which resolves or makes
    substantial progress on a long-standing open problem. These problems span
    high-dimensional geometry, coding theory, arithmetic circuit complexity,
    group theory, operator algebras, quantum complexity, lattice cryptography
    and extremal combinatorics.
  \end{quote}

  \begin{quote}
    The model formalized each argument in a Lean certificate.
  \end{quote}

  (OpenAI, Ten advances in mathematics and theoretical computer science\footnote{\url{https://openai.com/index/ten-advances-in-mathematics}})
\end{frame}

\begin{frame}{Formal Verification in the AI Era}
  \begin{itemize}
    \item From vibe coding to ``veri-coding''
      \begin{itemize}
        \item AIs write programs
        \item AIs write proofs
        \item Proofs are machine-checkable
      \end{itemize}
    \item Humans should decide and understand the proven properties
  \end{itemize}
\end{frame}

\begin{frame}{Curry-Howard Correspondence}
  \begin{itemize}
    \item The core mechanism behind proof assistants based on dependent type theory
      \begin{itemize}
        \item Proposition = Type
        \item Proof = Term (Program)
        \item Validating a proof = Type checking a term
      \end{itemize}
  \end{itemize}
\end{frame}

\begin{frame}{Today's Topics}
  \begin{itemize}
    \item Functions and inductive types
    \item Polymorphism
    \item Dependent types
    \item Curry-Howard correspondence
  \end{itemize}
\end{frame}

\begin{frame}{Live Coding}
  \begin{itemize}
    \item Material at \url{https://github.com/plrg-unist/lean-intro}
    \item Use VS Code or Lean Playground (\url{https://live.lean-lang.org/})
  \end{itemize}
\end{frame}

\begin{frame}{How to Verify Programs}
  \begin{itemize}
    \item Option 1: Write a program in a proof assistant and prove its properties
    \item Option 2: Embed a program written in another language into a proof
      assistant and prove its properties
  \end{itemize}
\end{frame}

\begin{frame}{More Resources}
  \begin{itemize}
    \item Software Foundations: \url{https://softwarefoundations.cis.upenn.edu/}
    \item The Hitchhiker's Guide to Logical Verification
  \end{itemize}
\end{frame}

\begin{frame}{Part 2 Preview}
  \begin{itemize}
    \item How can we state properties of programs?
    \item Can we have proof techniques that can be applied to a wide range of programs?
  \end{itemize}
\end{frame}

\end{document}
